% 382_Myon_g2_Stand2025_En_ch.tex
% Johann Pascher, 29 Sept. 2026
% Muon g-2: status 2025 and assessment of the FFGFT statements (successor document)

\providecommand{\KM}{\textbf{[K]}}
\providecommand{\SM}{\textbf{[S]}}

\chapter*{Muon \texorpdfstring{$g-2$}{g-2}: Status 2025}
\addcontentsline{toc}{chapter}{Muon g-2: Status 2025}

\begin{abstract}
Several early documents of the corpus refer to the deviation between the measured
anomalous magnetic moment of the muon and the Standard Model prediction as it stood in
2021: $\Delta a_\mu = 251(59)\times10^{-11}$, about $4.2\,\sigma$. With the final result
of the Fermilab experiment (June 2025) and the 2025 White Paper of the Muon $g-2$ Theory
Initiative, this deviation has shrunk to $38(63)\times10^{-11}$, about $0.6\,\sigma$
\KM. Statements that cite a T0 contribution of $251\times10^{-11}$ as the explanation of
the anomaly are therefore outdated; such a fixed contribution would today lie about
$3.4\,\sigma$ above the difference \KM.

The load-bearing $g-2$ line of the corpus (Doc.~018 Rev.~12, Doc.~158, A138) is not
affected. It rests on \emph{ratios}: $\Delta a(\tau{-}\mu)/\Delta a(\mu{-}e)=f^{1/3}-1$
and the bridge $\Delta a(\tau{-}e)/\Delta a(\mu{-}e)=(144/125)\,m_\tau/m_\mu$. In these,
the projection factor and the SI conversion cancel. It obtains absolute values by
anchoring the ratios to measured quantities; with the 2025 values this gives
$a_\tau=1.2811\times10^{-3}$ \KM. This anchored value is the better absolute value
compared with the direct formula, and it depends only on the measured values of $a_e$
and $a_\mu$, not on the Standard Model calculation, whose status has changed.
\end{abstract}

% ============================================================
\section*{Status markers}
% ============================================================
\begin{center}
\begin{tabular}{@{}lp{10cm}@{}}
\textbf{[K]} & Numerically confirmed (verification script). \\[3pt]
\textbf{[S]} & Speculative: a hint or conjecture, not yet derived or checked in the corpus. \\
\end{tabular}
\end{center}
The experimental and theoretical input values are taken from the sources listed in the
bibliography; they are not derived here.

% ============================================================
\section{Starting point in the corpus}
% ============================================================
The corpus contains two kinds of statements on the lepton $g-2$, which have to be kept
apart.

\textbf{(a) Explaining the anomaly by a T0 contribution.} Early documents posit a
contribution $\Delta a_\mu^{(\mathrm{T0})}=251\times10^{-11}$, or describe the muon
anomaly as explained by the $\xi$ term, or use it as a benchmark: Doc.~019, 030, 033,
049, 053, 059, 070, 081, 095, and in the A series A138 (comparison with the ``muon
discrepancy''). These statements depend directly on the size of the deviation between
measurement and Standard Model.

\textbf{(b) The ratio line.} Doc.~018 (Rev.~12), Doc.~158 and A138 describe the leptons
by fractal winding exponents ($5/3$ for the muon, $4/3$ for the tau) and put the ratios
at the centre: ratios of $g-2$ differences are structurally exact, because the
projection factor, the fractal correction and the SI conversion cancel. Absolute values
follow only by anchoring to a measured quantity -- the same division of labour as for
the masses, where ratios are fixed geometrically and one empirical anchor sets the scale
(A138, Doc.~351). This line compares with measured values, not with the difference to
the Standard Model.

% ============================================================
\section{Status 2025}
% ============================================================
\begin{itemize}
\item \textbf{Measurement.} The Fermilab experiment E989 published its final result in
  June 2025: $a_\mu=116\,592\,070.5(14.8)\times10^{-11}$, a precision of $127$~ppb
  \cite{FNAL2025}. Together with E821 (Brookhaven) this gives the world average
  $a_\mu^{\mathrm{exp}}=116\,592\,071.5(14.5)\times10^{-11}$ \cite{WP25}.
\item \textbf{Theory.} The 2025 White Paper of the Muon $g-2$ Theory Initiative gives
  $a_\mu^{\mathrm{SM}}=116\,592\,033(62)\times10^{-11}$ \cite{WP25}. The essential
  change compared with 2020 is the hadronic vacuum polarisation: it is now determined
  from lattice QCD instead of $e^+e^-$ cross sections.
\item \textbf{Difference.} $a_\mu^{\mathrm{exp}}-a_\mu^{\mathrm{SM}}=38(63)\times10^{-11}$,
  about $0.6\,\sigma$ \KM. In 2021 it was $251(59)\times10^{-11}$, about $4.2\,\sigma$
  \KM.
\item \textbf{Open within the theory.} Why the data-driven determination of the
  hadronic vacuum polarisation differs from the lattice determination is not yet
  understood \cite{CERNCourier2025}. The anomaly is therefore regarded as strongly
  weakened but not definitively settled.
\end{itemize}
What has changed is thus the Standard Model calculation, not the measured value: the
world average of $a_\mu$ has moved by only about $10\times10^{-11}$ since 2021.

% ============================================================
\section{Consequences for statements of type (a)}
% ============================================================
A fixed T0 contribution of $251\times10^{-11}$ was, in 2021, exactly as large as the
deviation it was meant to explain. Today it lies
\[
  \frac{251-38}{63}\approx3.4\,\sigma
\]
above the measured difference \KM. The statement ``the $\xi$ term explains the muon
anomaly'' is therefore outdated. For an additional contribution beyond the Standard
Model, the current status leaves a range of about $-88$ to $+164\times10^{-11}$ at two
standard deviations \KM. A T0 contribution of this size is not excluded, but it is not
derived in the corpus \SM.

The documents listed under (a) in Section~1 are affected. They carry a note after the
abstract that refers to this document; their text is otherwise unchanged. The
correction register records the change as R133 (Doc.~190).

% ============================================================
\section{The ratio line is not affected}
% ============================================================
The core of Doc.~018, Doc.~158 and A138 is two ratio statements \KM:
\[
  \frac{\Delta a(\tau{-}\mu)}{\Delta a(\mu{-}e)} = f^{1/3}-1 = 18.57,
  \qquad
  \frac{\Delta a(\tau{-}e)}{\Delta a(\mu{-}e)} = \frac{144}{125}\,\frac{m_\tau}{m_\mu} = 19.37 .
\]
The first follows from the exponents $5/3$ and $4/3$ alone. The second replaces
$f^{1/3}=19.57$ by the ratio of the mass formulas and thereby links $g-2$ to the lepton
masses; $f$, $k_\mathrm{geom}$ and $K_\mathrm{frak}$ cancel.

\textbf{From ratio to absolute value.} Anchoring the bridge to the measured values of
$a_e$ and $a_\mu$ gives the absolute value for the tau:
\[
  a_\tau = a_e + \frac{144}{125}\,\frac{m_\tau}{m_\mu}\,(a_\mu-a_e) = 1.2811\times10^{-3}
\]
with the 2025 values \KM; Doc.~018 and A138 quote $1.282\times10^{-3}$ with a rounded
$a_e$, Doc.~158 $1.281\times10^{-3}$. The direct absolute formula
$a_\tau=a_e+4\pi/f^{4/3}$ gives $1.2454\times10^{-3}$ instead \KM. The anchored value is
the better one, because it contains only the exact ratio and measured quantities; this
is why Doc.~158 and A138 use it as the prediction.

\textbf{Why the anomaly plays no role here.} The anchor is the measured difference
$a_\mu-a_e=6.2685\times10^{-6}$. It has changed by only about $10^{-10}$ between 2021 and
2025 \KM. The anomaly was a statement about the Standard Model calculation; the ratio
line does not use that calculation. The tau prediction therefore remains unchanged. It
lies $1.04\times10^{-4}$ above the Standard Model prediction
$a_\tau^{\mathrm{SM}}=1.1772\times10^{-3}$ \KM. The current experimental bounds on
$a_\tau$ are about 600 times wider than this difference \KM; the decision awaits a
measurement of $a_\tau$ \SM.

% ============================================================
\section{The direct absolute formulas}
% ============================================================
Besides the ratio line, Doc.~018 contains direct absolute formulas from the winding
geometry, corrected by $K_\mathrm{frak}^{3/2}$. With the values of Doc.~018 one obtains
\KM:
\begin{align*}
  a_e^{(\mathrm{corr})} &= 1.1598\times10^{-3}\quad(0.014\,\%\ \text{from the measured value}),\\
  a_\mu^{(\mathrm{corr})} &= 1.1642\times10^{-3}\quad(-0.15\,\%\ \text{from the world average}).
\end{align*}
These values also compare with measured values and do not depend on the anomaly. They
work at the level of $10^{-4}$ to $10^{-3}$, however, and are not a precision comparison
where the anomaly was discussed. A numerical example shows why the corpus takes the
route via ratios for absolute values: the directly computed difference
$\Delta a(\mu{-}e)=4\pi/f^{5/3}=4.373\times10^{-6}$ lies about $30\,\%$ below the measured
one \KM, while the ratio of the differences follows exactly from the exponents. The
anchored route avoids this uncertainty.

\textbf{A138.} The comparison there, that the tau deviation is ``seven times larger than
the muon discrepancy'', has become moot, since there is no longer a significant muon
discrepancy. The tau statement itself is independent of this.

% ============================================================
\section{Result}
% ============================================================
With the 2025 status, the muon $g-2$ anomaly of 2021 is no longer significant. FFGFT
statements that use it as evidence or as the target of a T0 contribution are outdated
and are marked in the affected documents. The load-bearing ratio line is not affected:
its ratios are structurally exact, it obtains its absolute values by anchoring to
measured quantities, and the measured values have not changed. Its tau prediction
$a_\tau\approx1.281\times10^{-3}$ remains the actual test statement. Whether the FFGFT
additionally yields a contribution to the muon within today's allowed range of about
$10^{-9}$ is open \SM.

\medskip\noindent Verification script: \texttt{2/python/Dok382\_Skripte/pruef\_382\_myon\_g2\_stand.py}, 21/21.

\begin{thebibliography}{9}
\bibitem{FNAL2025} Muon $g-2$ Collaboration (D.~P.~Aguillard et al.),
  \emph{Measurement of the Positive Muon Anomalous Magnetic Moment to 127~ppb},
  Phys.\ Rev.\ Lett.\ 135 (2025) 101802, arXiv:2506.03069.
\bibitem{WP25} R.~Aliberti et al. (Muon $g-2$ Theory Initiative),
  \emph{The anomalous magnetic moment of the muon in the Standard Model: an update},
  Phys.\ Rep.\ 1143 (2025) 1--158, arXiv:2505.21476.
\bibitem{CERNCourier2025} CERN Courier, \emph{Fermilab's final word on muon g-2} (2025).
\end{thebibliography}
